\subsection{Deciding at arrival}
\label{sec:arrival}
A policy need not wait for a split: it can rank each block when it is created, against the blocks already
placed. The decision that matters is the same one: where a block with no history ranks against a block
seen once. Every block in a static analysis gets a belief that reads nothing after its own creation;
beliefs are placed in one capacity-constrained assignment. Ranking new blocks among themselves (by
arrival order, by a learned score from session context, or at random) moves the result by at most 5.3
points; moving them ahead of or behind equal-count visible blocks moves it by 14--29 points. An \emph{estimated
prior} replaces that tie rule: at a pseudo-split halfway through the prefix, blocks born after it are
reused before the real split more than blocks already seen once (0.52 against 0.30 on conversation), so
new blocks are placed between count 2 and count 3. It captures 47.0\% and 51.7\% with tier-only actions
(58.2\% and 61.1\% with all actions), and nothing on the synthetic control with tier-only actions. On Alibaba's four Qwen
serving traces~\cite{wang2025kvcachewild} (16-token blocks) it captures 54.3\%, 34.2\% and 45.9\% on three
and 0.0\% on the fourth (thinking), where the prefix rates seen-once blocks above new ones. The direction
the prefix gives is a regularity of these workloads, not a law, and the estimate has to be checked per
workload.
